\documentclass[11pt,a4paper]{article}

% --- PAQUETES BÁSICOS Y TIPOGRAFÍA ---
\usepackage[utf8]{inputenc}
\usepackage[spanish,es-tabla]{babel}
\usepackage[margin=2.2cm]{geometry}
\usepackage{amsmath,amssymb}
\usepackage{booktabs}
\usepackage{array}
\usepackage{graphicx}
\usepackage{microtype}
\usepackage{listings}
\usepackage{xcolor}
\usepackage{hyperref}
\usepackage{fancyhdr}
\usepackage{titlesec}

% --- CONFIGURACIÓN DE ENLACES Y COLORES ---
\definecolor{codegreen}{rgb}{0,0.6,0}
\definecolor{codegray}{rgb}{0.5,0.5,0.5}
\definecolor{codepurple}{rgb}{0.58,0,0.82}
\definecolor{backcolour}{rgb}{0.96,0.96,0.96}
\definecolor{navyblue}{rgb}{0.0, 0.0, 0.5}

\hypersetup{
    colorlinks=true,
    linkcolor=navyblue,
    citecolor=navyblue,
    urlcolor=navyblue
}

% --- CONFIGURACIÓN DE LISTINGS (C / OPENMP) ---
\lstdefinestyle{cstyle}{
    backgroundcolor=\color{backcolour},   
    commentstyle=\color{codegreen},
    keywordstyle=\color{navyblue}\bfseries,
    numberstyle=\tiny\color{codegray},
    stringstyle=\color{codepurple},
    basicstyle=\ttfamily\footnotesize,
    breakatwhitespace=false,         
    breaklines=true,                 
    captionpos=b,                    
    keepspaces=true,                 
    numbers=left,                    
    numbersep=5pt,                  
    showspaces=false,                
    showstringspaces=false,
    showtabs=false,                  
    tabsize=2,
    frame=single
}
\lstset{style=cstyle}

% --- ENCABEZADOS Y PIES DE PÁGINA ---
\pagestyle{fancy}
\fancyhf{}
\fancyhead[L]{\small Laboratorio PCAM + OpenMP: Simulación N-Body 2D}
\fancyhead[R]{\small Entrega Oficial}
\fancyfoot[C]{\thepage}

% ==============================================================================
% DATOS DEL INFORME
% ==============================================================================
\title{
    \vspace{-1.5cm}
    \textbf{\Large Laboratorio de Computación Paralela: Metodología PCAM y OpenMP}\\[0.2cm]
    \textbf{\large Simulación Gravitacional N-Body 2D}
}

\author{
    \textbf{Jonathan Zacarias} (Base, Corrección y Ensamblado) \\
    \textbf{Hugo Barillas} (Paralelo y Experimentos) \\
    \textbf{June Herrera} (Optimización y Concurrencia)
}
\date{Octubre 2026}

% ==============================================================================
\begin{document}

\maketitle

\begin{abstract}
En este laboratorio se presenta el diseño, desarrollo, paralelización y optimización de una simulación gravitatoria bidimensional de $N$-cuerpos ($N$-Body 2D) bajo la metodología sistemática \textbf{PCAM} (\textit{Partitioning, Communication, Agglomeration, Mapping}) y el estándar de memoria compartida \textbf{OpenMP} en lenguaje C. Se formula la línea base secuencial con complejidad algorítmica $\mathcal{O}(N^2)$ y se establecen mecanismos rigurosos de validación numérica mediante la definición de un \textit{checksum} acumulado determinista con semilla fija ($42$). El presente documento detalla la fase de particionamiento, la estrategia de desacoplamiento temporal, la corrección numérica de la simulación y las respuestas a las interrogantes de escalabilidad planteadas.
\end{abstract}

\vspace{0.3cm}

% ==============================================================================
\section{Introducción y Línea Base Secuencial (Baseline)}
% Responsable: Jonathan Zacarias
% ==============================================================================

El problema de simulación de $N$-cuerpos (\textit{N-Body Problem}) modela la evolución dinámica en el tiempo de un conjunto de $N$ partículas que interactúan mutuamente según la Ley de Gravitación Universal. En un espacio euclidiano bidimensional, cada cuerpo $i \in \{0, \dots, N-1\}$ está caracterizado por su posición $\mathbf{r}_i = (x_i, y_i)$, velocidad $\mathbf{v}_i = (v_{x,i}, v_{y,i})$ y masa $m_i > 0$.

La fuerza gravitatoria neta $\mathbf{F}_i$ que experimenta la partícula $i$ debido a la atracción de las restantes $N-1$ partículas viene dada por:
\begin{equation}
    \mathbf{F}_i = \sum_{\substack{j=0 \\ j \neq i}}^{N-1} G \frac{m_i m_j}{\left( \|\mathbf{r}_j - \mathbf{r}_i\|^2 + \epsilon^2 \right)^{3/2}} (\mathbf{r}_j - \mathbf{r}_i)
    \label{eq:fuerza}
\end{equation}
donde $G = 1.0$ representa la constante gravitacional normalizada y $\epsilon^2 = 10^{-9}$ es un factor de suavizado (\textit{softening parameter}) indispensable para evitar singularidades y divisiones por cero cuando la distancia euclidiana entre dos partículas tiende a cero.

A partir de la aceleración resultante $\mathbf{a}_i = \mathbf{F}_i / m_i$, la evolución cinemática se resuelve discretizando el tiempo mediante un integrador de Euler con paso temporal $\Delta t = 0.01\text{ s}$:
\begin{align}
    \mathbf{v}_i(t + \Delta t) &= \mathbf{v}_i(t) + \mathbf{a}_i(t) \cdot \Delta t \\
    \mathbf{r}_i(t + \Delta t) &= \mathbf{r}_i(t) + \mathbf{v}_i(t + \Delta t) \cdot \Delta t
\end{align}

\subsection{Complejidad Algorítmica y Desacoplamiento en Dos Fases}
El cálculo directo requiere evaluar $N(N - 1) = \mathcal{O}(N^2)$ interacciones por cada paso de tiempo, mientras que la integración física demanda únicamente $\mathcal{O}(N)$ operaciones. Para la configuración de referencia ($N = 5000$ cuerpos, $10$ pasos), se computan $2.5 \times 10^7$ interacciones por paso, consumiendo más del 99.5\% del tiempo en el cálculo de fuerzas.

La implementación secuencial de referencia (\texttt{nbody\_secuencial.c} / \texttt{src/nbody\_seq.c}) estructura cada paso en \textbf{dos fases desacopladas}:
\begin{enumerate}
    \item \textbf{Fase 1 (Cómputo de Fuerzas):} Se calculan y acumulan las componentes de fuerza en vectores temporales \texttt{fx[i]} y \texttt{fy[i]} leyendo exclusivamente las posiciones de los cuerpos.
    \item \textbf{Fase 2 (Actualización Cinemática):} Se actualizan velocidades y posiciones en un bucle separado de orden $\mathcal{O}(N)$.
\end{enumerate}
Este desacoplamiento garantiza que ninguna partícula altere su estado durante el cómputo de fuerzas de un mismo paso, eliminando dependencias cruzadas de lectura y escritura para la posterior paralelización.

% ==============================================================================
\section{Aplicación de la Metodología PCAM}
% ==============================================================================

\subsection{Partition (Particionamiento) - [Responsable: Jonathan Zacarias]}
La fase de partición descompone el cómputo en tareas concurrentes elementales de grano fino. En la simulación $N$-Body se identificaron dos estrategias principales:

\begin{enumerate}
    \item \textbf{Partición por Partícula / Grano Medio ($\mathcal{O}(N)$ tareas):}
    Se define la tarea $T_i$ como el cálculo integral de la fuerza neta ejercida sobre el cuerpo $i$. Cada tarea $T_i$ requiere acceso de lectura sobre el arreglo global de posiciones de todas las partículas, pero escribe \textit{únicamente} en sus acumuladores privados \texttt{fx[i]} y \texttt{fy[i]}. Esta partición expone $N$ tareas estrictamente independientes y libres de condiciones de carrera por diseño (\textit{race-free}), constituyendo la base de \texttt{nbody\_paralelo.c}.
    
    \item \textbf{Partición por Pares de Interacción / Grano Fino ($\mathcal{O}(N^2 / 2)$ tareas):}
    Aprovecha la Tercera Ley de Newton ($\mathbf{F}_{ji} = -\mathbf{F}_{ij}$). Cada tarea $T_{i,j}$ (con $j > i$) calcula una interacción y actualiza simultáneamente los cuerpos $i$ y $j$. Si bien reduce el número de operaciones a la mitad, introduce escrituras concurrentes sobre acumuladores compartidos, requiriendo sincronización o reducción en variables locales por hilo (\texttt{nbody\_optimized.c}).
\end{enumerate}

\subsection{Communication y Race Conditions - [Responsable: June Herrera]}
%% --- SECCIÓN PARA JUNE HERRERA ---
\textit{[Aquí June Herrera describe el análisis de comunicación en memoria compartida, el acceso al bus de memoria/caché y la resolución de condiciones de carrera en nbody\_optimized.c mediante acumuladores privados o directivas OpenMP atomic/critical].}

\subsection{Agglomeration y Mapping - [Responsable: Hugo Barillas]}
%% --- SECCIÓN PARA HUGO BARILLAS ---
\textbf{Agglomeration.} Asignar cada tarea $T_i$ por separado haría que el costo de repartirlas pese demasiado, así que las iteraciones consecutivas del ciclo sobre $i$ se agrupan en bloques de tamaño \texttt{chunk}. Se comparan tres granularidades con \texttt{static} y 8 hilos: \texttt{chunk} por defecto (un bloque de $N/p = 625$ cuerpos por hilo), 64 y 8. Un bloque más pequeño reparte el trabajo de forma más fina, pero implica más bloques que administrar y más bordes compartidos en \texttt{fx[]}/\texttt{fy[]} (\textit{false sharing}). Como cada bloque de 8 cuerpos ya contiene $8 \times 4999 \approx 4\times10^4$ interacciones, el costo de administrar bloques es despreciable. Con \texttt{static}, el chunk cambia el tiempo menos que la variación entre corridas; con \texttt{dynamic}, en cambio, el chunk 8 superó al 64 (Tabla~\ref{tab:experimentos} y pregunta 3). Además, las dos fases de cada paso (fuerzas e integración) se aglomeran dentro de \textbf{una sola región paralela}, de modo que el equipo de hilos se crea una sola vez y no en cada paso.

\textbf{Mapping.} Los bloques se asignan a los hilos con \texttt{schedule(runtime)}, y la política se elige por línea de comandos sin recompilar:
\begin{itemize}
    \item \texttt{static}: asignación fija y sin costo durante la ejecución. A priori es la mejor opción, porque cada $T_i$ cuesta exactamente $N-1$ interacciones.
    \item \texttt{dynamic}: cada hilo pide un bloque nuevo al terminar el anterior. Cuesta una operación de sincronización por bloque, pero compensa diferencias de velocidad entre hilos que no vienen del algoritmo sino del hardware: núcleos con distinta frecuencia turbo, dos hilos de \textit{hyperthreading} en un mismo núcleo y procesos del sistema operativo.
    \item \texttt{guided}: empieza con bloques grandes y los reduce progresivamente.
\end{itemize}
La fase de integración, de costo $\mathcal{O}(1)$ por cuerpo, siempre usa \texttt{schedule(static)}.

% ==============================================================================
\section{Decisiones de Diseño con OpenMP}
% Responsables: Hugo Barillas y June Herrera
% ==============================================================================
%% --- SECCIÓN PARA HUGO BARILLAS Y JUNE HERRERA ---
\textbf{Versión paralela directa (\texttt{nbody\_paralelo.c}).}
\begin{itemize}
    \item \textbf{Una sola región \texttt{\#pragma omp parallel}} cubre todo el ciclo temporal, con dos \texttt{\#pragma omp for} por paso, en lugar de abrir un \texttt{parallel for} en cada paso. Así se evita pagar el \textit{fork-join} en cada paso.
    \item \textbf{\texttt{default(none)}} obliga a declarar el alcance de cada variable. Son \texttt{shared}: \texttt{bodies}, \texttt{fx}, \texttt{fy}, \texttt{n}, \texttt{steps} y \texttt{threads\_used}. Son privadas, por estar declaradas dentro de la región: \texttt{step}, \texttt{i}, \texttt{j}, los acumuladores \texttt{fxi} y \texttt{fyi}, y los temporales (\texttt{dx}, \texttt{dy}, \texttt{dist}, \ldots). No hace falta \texttt{firstprivate}.
    \item \textbf{Sin condiciones de carrera.} En la fase de fuerzas, la iteración $i$ solo escribe \texttt{fx[i]} y \texttt{fy[i]}, y \texttt{bodies[]} es de solo lectura. En la integración, la iteración $i$ solo modifica \texttt{bodies[i]}. Por eso no se necesitan \texttt{atomic}, \texttt{critical} ni \texttt{reduction}.
    \item \textbf{Barreras implícitas} al final de cada \texttt{omp for} (sin \texttt{nowait}). La primera asegura que todas las fuerzas estén calculadas antes de mover cualquier cuerpo. La segunda evita que un hilo empiece el paso siguiente con posiciones a medio actualizar.
    \item La directiva \texttt{\#pragma omp single} registra, con \texttt{omp\_get\_num\_threads()}, los hilos que el runtime entregó realmente; se imprimen junto con los solicitados.
    \item El tiempo se mide con \texttt{omp\_get\_wtime()} y abarca solo la simulación. Como cada $T_i$ suma sobre $j$ en el mismo orden que la versión secuencial, el checksum es idéntico bit a bit en todas las configuraciones (verificado con 17 dígitos significativos).
\end{itemize}

% ==============================================================================
\section{Corrección Numérica y Validación de Resultados}
% Responsable: Jonathan Zacarias
% ==============================================================================

Para verificar la corrección numérica de las implementaciones paralelas frente al baseline secuencial, se formuló el \textbf{Checksum acumulado} $C(t)$ al finalizar los 10 pasos temporales:
\begin{equation}
    C = \sum_{i=0}^{N-1} \left( x_i + y_i + v_{x,i} + v_{y,i} \right)
    \label{eq:checksum}
\end{equation}

\subsection{Análisis de Asociatividad en Aritmética Flotante (IEEE-754)}
En la aritmética de coma flotante de doble precisión (\texttt{double}), la adición no es asociativa: $(a + b) + c \neq a + (b + c)$.
\begin{itemize}
    \item En la versión paralela (\texttt{nbody\_paralelo.c}), como cada hilo ejecuta el bucle interno $j$ en el orden exacto $0 \dots N-1$, el Checksum resultante debe ser \textbf{estrictamente idéntico} al secuencial ($\Delta C = 0$).
    \item En la versión optimizada (\texttt{nbody\_optimized.c}), las reducciones parciales alteran el orden cronológico de suma, tolerándose un error numérico relativo:
    \begin{equation}
        \epsilon_{\text{rel}} = \frac{|C_{\text{opt}} - C_{\text{seq}}|}{|C_{\text{seq}}|} < 10^{-11}
    \end{equation}
    Cualquier desviación superior a $10^{-7}$ evidencia una condición de carrera activa.
\end{itemize}

\begin{table}[h!]
\centering
\caption{Validación de corrección numérica con semilla $42$ (10 pasos temporales)}
\label{tab:correccion}
\begin{tabular}{@{}ccccc@{}}
\toprule
\textbf{Cuerpos ($N$)} & \textbf{Checksum Secuencial} & \textbf{Checksum Paralelo} & \textbf{Error Relativo} & \textbf{Estado} \\ \midrule
100 & Baseline Ref. & Idéntico & $\Delta = 0$ & Validado \\
1000 & Baseline Ref. & Idéntico & $\Delta = 0$ & Validado \\
3000 & Baseline Ref. & Idéntico & $\Delta = 0$ & Validado \\
5000 & Baseline Ref. & Idéntico & $\Delta = 0$ & Validado \\ \bottomrule
\end{tabular}
\end{table}

% ==============================================================================
\section{Resultados Experimentales, Speedup y Eficiencia}
% Responsable: Hugo Barillas
% ==============================================================================
%% --- SECCIÓN PARA HUGO BARILLAS ---
% Tabla generada por scripts/analyze_results.py a partir de results.csv.
Cada configuración se ejecutó 5 veces ($\geq 3$ exigidas), con $N = 5000$, 10 pasos y semilla 42, en un Intel Core i7-1165G7 (4 núcleos físicos, 8 lógicos) conectado a la corriente. Las repeticiones se \textbf{intercalan}: cada ronda recorre todas las configuraciones, para que la deriva térmica y de frecuencia de la laptop afecte a todas por igual. Antes de la primera ronda se hace una corrida de calentamiento que se descarta. La Tabla~\ref{tab:experimentos} reporta el promedio $\bar{T}$, la desviación estándar $\sigma$, $S(p) = T_1/T_p$, $E(p) = S(p)/p$ y la fracción serial experimental de Karp-Flatt, $e_{KF} = \frac{1/S - 1/p}{1 - 1/p}$, con $T_1$ el promedio de la versión secuencial. Todas las corridas reproducen exactamente el checksum secuencial (error relativo $0$).

\begin{table}[h!]
\centering
\caption{Tiempos promedio (5 repeticiones), speedup y eficiencia. $N=5000$, 10 pasos, $T_1 = 1.788$ s}
\label{tab:experimentos}
\small
\begin{tabular}{@{}lcccrrrrr@{}}
\toprule
\textbf{Versión} & \textbf{$p$} & \textbf{Schedule} & \textbf{Chunk} & \textbf{$\bar{T}$ (s)} & \textbf{$\sigma$ (s)} & \textbf{$S(p)$} & \textbf{$E(p)$} & \textbf{$e_{KF}$} \\ \midrule
nbody\_secuencial & 1 & sequential & default & 1.788 & 0.245 & 1.00 & 1.00 & -- \\
nbody\_paralelo & 1 & static & default & 1.781 & 0.168 & 1.00 & 1.00 & -- \\
nbody\_paralelo & 2 & static & default & 1.082 & 0.087 & 1.65 & 0.83 & 0.211 \\
nbody\_paralelo & 4 & static & default & 0.748 & 0.054 & 2.39 & 0.60 & 0.225 \\
nbody\_paralelo & 8 & static & default & 0.547 & 0.025 & 3.27 & 0.41 & 0.207 \\
nbody\_paralelo & 8 & static & 8 & 0.542 & 0.031 & 3.30 & 0.41 & 0.204 \\
nbody\_paralelo & 8 & static & 64 & 0.570 & 0.048 & 3.14 & 0.39 & 0.222 \\
nbody\_paralelo & 8 & dynamic & 8 & 0.474 & 0.028 & 3.77 & 0.47 & 0.160 \\
nbody\_paralelo & 8 & dynamic & 64 & 0.519 & 0.025 & 3.45 & 0.43 & 0.189 \\
nbody\_paralelo & 8 & guided & default & 0.537 & 0.049 & 3.33 & 0.42 & 0.200 \\
\bottomrule
\end{tabular}
\end{table}


\textbf{Escalabilidad.} Con \texttt{static}: $S(2)=1.65$, $S(4)=2.39$ y $S(8)=3.27$, es decir, $E = 0.83$, $0.60$ y $0.41$. La versión paralela con 1 hilo tarda lo mismo que la secuencial ($S = 1.00$), así que crear la región paralela y pasar por las barreras no cuesta nada medible. Además, $e_{KF}$ se mantiene casi constante ($\approx 0.21$) entre 2 y 8 hilos. Si la pérdida viniera de sincronización u overhead de OpenMP, $e_{KF}$ crecería con $p$. Por lo tanto, la pérdida es un costo aproximadamente fijo de usar más de un núcleo en este equipo: en una laptop los núcleos comparten el presupuesto de potencia y temperatura, y la frecuencia por núcleo baja con la carga. De 4 a 8 hilos se gana solo $1.37\times$, porque los hilos 5 a 8 son de \textit{hyperthreading} y comparten las unidades de punto flotante de núcleos ya ocupados.

\textbf{Schedules con 8 hilos.} El mejor resultado fue \texttt{dynamic} con chunk 8: $0.474$~s ($S = 3.77$), contra $0.542$--$0.570$~s de \texttt{static} y $0.537$~s de \texttt{guided}. La diferencia (13\,\%) es mayor que dos desviaciones estándar. La carga algorítmica es uniforme, pero los hilos no avanzan a la misma velocidad: con \texttt{static}, cada paso termina cuando acaba el hilo más lento, mientras que con \texttt{dynamic} los hilos más rápidos toman más bloques. Esto también explica el menor $e_{KF}$ de \texttt{dynamic} (0.16 frente a 0.21).

% ==============================================================================
\section{Comparación: Versión Paralela vs. Optimizada}
% Responsable: June Herrera
% ==============================================================================
%% --- SECCIÓN PARA JUNE HERRERA ---
\textit{[Aquí June Herrera compara el rendimiento de la paralelización O(N^2) frente a la reducción simétrica O(N^2 / 2) con pairs y analiza el balance de carga].}

% ==============================================================================
\section{Conclusiones}
% Responsable: Todo el equipo
% ==============================================================================
%% --- CONCLUSIONES CONJUNTAS ---
\textit{[Conclusiones generales sobre la escalabilidad, el impacto de la sincronización en OpenMP y la metodología PCAM].}

\newpage
% ==============================================================================
\appendix
\section{Anexo: Respuestas a Preguntas de Análisis}
% ==============================================================================

\subsection*{Pregunta 1: Influencia del tamaño del problema ($N$) en la relación cómputo/comunicación y en el Speedup [Responsable: Jonathan Zacarias]}

\textbf{Respuesta:}  
En la simulación N-Body directa, la complejidad computacional del cálculo de fuerzas es cuadrática $\mathcal{O}(N^2)$, mientras que la sobrecarga de comunicación (accesos al bus de memoria compartida, barreras de sincronización y creación de hilos en el modelo \textit{fork-join}) es de orden $\mathcal{O}(N)$ o proporcional al número de hilos $\mathcal{O}(p)$.

Para valores pequeños de $N$ ($N \le 500$), el tiempo de cálculo aritmético es mínimo. En este régimen, la sobrecarga impuesta por el runtime de OpenMP para inicializar los hilos y sincronizarlos al término de cada paso domina sobre el cómputo útil, resultando en un \textit{Speedup} deficiente o incluso inferior a la unidad ($S < 1$, desaceleración neta).

Por el contrario, para valores grandes de $N$ ($N \ge 3000$ y $N = 5000$), el número de evaluaciones de fuerza crece hasta $2.5 \times 10^7$ operaciones por paso. La relación entre cómputo útil y sobrecarga se incrementa cuadráticamente. Conforme a la \textbf{Ley de Amdahl} y la \textbf{Ley de Gustafson-Barsis}, la fracción paralelizable del código supera el 99.5\%, amortizando completamente el \textit{overhead} de OpenMP y logrando un escalado casi lineal para 2, 4 y 8 hilos.

\vspace{0.5cm}

\subsection*{Pregunta 2: Impacto de la granularidad de la partición y el balance de carga [Responsable: Jonathan Zacarias]}

\textbf{Respuesta:}  
La granularidad de la partición define el tamaño de las tareas independientes y su grado de uniformidad computacional:

\begin{enumerate}
    \item \textbf{En la versión paralela (\texttt{nbody\_paralelo.c}):}
    Al particionar el bucle externo sobre las partículas $i$, cada tarea elemental ejecuta exactamente el mismo número de iteraciones ($N-1$ evaluaciones). La carga es estrictamente homogénea. Por ende, la política estática (\texttt{schedule(static)}) divide equitativamente las partículas en bloques contiguos de tamaño $N/p$, logrando un balance de carga óptimo sin incurrir en la sobrecarga de gestión dinámica de colas.
    
    \item \textbf{En la versión optimizada (\texttt{nbody\_optimized.c}):}
    Al explotar la simetría mediante pares $(i, j)$ con $j > i$, el cuerpo $i=0$ evalúa $N-1$ interacciones, mientras que el cuerpo $i=N-2$ evalúa una sola interacción. Esto genera una matriz triangular superior con una distribución de carga intrínsecamente desbalanceada. Si se utilizara una partición estática ingenua, los primeros hilos soportarían la mayor parte del trabajo mientras los últimos quedarían ociosos esperando en la barrera. En este escenario, una granularidad más fina combinada con planificaciones dinámicas (\texttt{schedule(dynamic, chunk)}) o guiadas (\texttt{schedule(guided)}) resulta imprescindible para mitigar el desbalance computacional.
\end{enumerate}

\vspace{0.5cm}

\subsection*{Preguntas 3, 4, 7 y 8 [Responsable: Hugo Barillas]}
%% --- RESPUESTAS HUGO BARILLAS ---
\textbf{Pregunta 3. ¿Qué efecto tuvo el chunk size sobre el rendimiento? Relaciónelo con Agglomeration.}
Con \texttt{static} y 8 hilos, el chunk por defecto (625), 8 y 64 dieron 0.547, 0.542 y 0.570~s. Las diferencias son menores que la desviación entre corridas (0.03--0.05~s), así que ahí el chunk no influye de forma medible: aun con chunk 8, cada bloque agrupa unas $4\times10^4$ interacciones y el costo de administrarlo nunca domina. Con \texttt{dynamic} el chunk sí importó: chunk 8 tardó 0.474~s y chunk 64, 0.519~s ($\sigma \approx 0.03$). Con chunk 64 solo hay $5000/64 \approx 78$ bloques para 8 hilos; al final de cada paso, un hilo puede quedarse con un último bloque de 64 cuerpos mientras los demás esperan en la barrera. Con chunk 8, ese residuo es 8 veces menor. En términos de Agglomeration, conviene agrupar lo suficiente para que el reparto no domine, pero dejar bloques lo bastante finos para que el balance dinámico funcione. Con este $N$, un chunk de 8 cumple ambas condiciones.

\textbf{Pregunta 4. Compare static, dynamic y guided. ¿Cuál funcionó mejor y qué evidencia lo respalda?}
\texttt{dynamic} (chunk 8) fue el mejor: $0.474 \pm 0.028$~s, $S = 3.77$, frente a \texttt{static} ($0.542$--$0.570$~s, $S \le 3.30$) y \texttt{guided} ($0.537 \pm 0.049$~s, $S = 3.33$). La diferencia supera dos desviaciones estándar, y la tendencia se repitió en una corrida previa independiente de 3 repeticiones. Con \texttt{dynamic}, $e_{KF}$ baja de $\approx 0.21$ a $0.16$, lo que indica que \texttt{static} perdía tiempo por desbalance. Aunque cada cuerpo cuesta lo mismo ($N-1$ interacciones), los hilos no corren a la misma velocidad: dos hilos de hyperthreading comparten un núcleo, la frecuencia varía y el sistema operativo interrumpe. \texttt{static} no puede compensarlo y \texttt{dynamic} sí, con un costo de scheduling despreciable. \texttt{guided} empieza con bloques grandes que fijan la mayor parte del reparto, por lo que se comporta casi como \texttt{static}.

\textbf{Pregunta 7. Speedup y Efficiency para 2, 4 y 8 threads. ¿Dónde dejó de escalar linealmente?}
Con \texttt{static}: $S(2)=1.65$, $E(2)=0.83$; $S(4)=2.39$, $E(4)=0.60$; $S(8)=3.27$, $E(8)=0.41$. Con \texttt{dynamic} (chunk 8), $S(8)=3.77$ y $E(8)=0.47$. El escalamiento deja de ser lineal desde 2 hilos, no por overhead de OpenMP (paralelo con 1 hilo da $S=1.00$ y $e_{KF}$ no crece con $p$), sino porque cada núcleo adicional rinde menos que el primero en esta laptop (potencia y temperatura compartidas). La pendiente cae aún más de 4 a 8 hilos ($1.37\times$), al pasar de núcleos físicos a hilos de hyperthreading.

\textbf{Pregunta 8. Con 64 cores, ¿esperaría 8 veces el rendimiento obtenido con 8 threads?}
No. Si se toma la fracción no paralelizable medida, $e_{KF} \approx 0.16$--$0.21$, la ley de Amdahl da $S(64) = 1/(e + (1-e)/64) \approx 4.5$--$5.8$, apenas $1.4$--$1.5$ veces el $S(8)$ obtenido y muy lejos de 8 veces. Esa extrapolación es pesimista, porque buena parte de nuestro $e_{KF}$ viene de limitaciones de la laptop (potencia, frecuencia y hyperthreading) que un servidor con 64 núcleos físicos reduciría. Aun así, hay límites que no desaparecen. Con $N = 5000$ y 64 hilos cada hilo recibe unos 78 cuerpos por paso, así que las 20 barreras y el desbalance del último bloque pesan proporcionalmente más. Además, el ancho de banda de memoria y la caché compartida no crecen al mismo ritmo que los núcleos. Para acercarse a un escalamiento lineal con 64 núcleos habría que aumentar $N$ (escalamiento de Gustafson), de modo que el trabajo $\mathcal{O}(N^2)$ crezca más rápido que los costos fijos de sincronización.

\subsection*{Preguntas 5 y 6 [Responsable: June Herrera]}
%% --- RESPUESTAS JUNE HERRERA ---
\textit{[Espacio reservado para las respuestas asignadas a June Herrera].}

\end{document}
